\documentclass[11pt]{article}
\usepackage[T1]{fontenc}
\usepackage{lmodern}
\usepackage{amsmath,amssymb}
\usepackage[margin=1in]{geometry}
\usepackage[colorlinks=true,linkcolor=blue,urlcolor=blue,citecolor=blue]{hyperref}

\newcommand{\Md}{M_d}

\title{Literature status: direct-product stability of the $M_d$-approximation property}
\author{Run \texttt{fa\_banach\_001}, lane 2}
\date{11 August 2026}

\begin{document}
\maketitle

\paragraph{Status.}
\texttt{literature\_already\_answered}. This packet records an exact later
answer and does not claim a new mathematical result.

\paragraph{Source question.}
In Remark 4.2 on PDF page 7 of Vergara's 2023 paper
\cite{Vergara2023}, the author writes: ``We do not know if $M_d$-AP is
stable under direct products.'' The setting is discrete groups, with
$M_d(G)=X_d(G)^*$ and $M_d$-AP meaning that the constant function $1$ lies
in the $\sigma(M_d(G),X_d(G))$-closure of $\mathbb C[G]$.

\paragraph{Exact later answer.}
Vergara's 2025 follow-up \cite{Vergara2025} explicitly cites the source and
proves a stronger extension theorem. Its Theorem 1.3 states that, for a
normal subgroup $\Gamma\triangleleft G$ and $d\geq2$, if both $\Gamma$ and
$G/\Gamma$ have $M_d$-AP, then $G$ has $M_d$-AP. Corollary 1.4 on PDF page
3 states explicitly that $M_d$-AP is stable under direct products,
semidirect products, and free products. The proof on PDF page 10 obtains the
direct-product case from the semidirect-product case with trivial action.
Thus Remark 4.2 is answered affirmatively.

\paragraph{Independent consistency check.}
If $u_i\to1$ weak-* in $M_d(G)$ and $v_j\to1$ weak-* in $M_d(H)$ are
finitely supported, set $w_{i,j}(g,h)=u_i(g)v_j(h)$. Tensoring the defining
Hilbert-space factorizations gives
\[
 \|w_{i,j}\|_{M_d(G\times H)}
 \leq \|u_i\|_{M_d(G)}\,\|v_j\|_{M_d(H)}.
\]
Both approximating nets are norm bounded. The product net converges pointwise
to $1$ and therefore converges on $\mathbb C[G\times H]$; its uniform norm
bound and the density of this subspace in $X_d(G\times H)$ upgrade the
convergence to weak-* convergence. This verifies the direct-product statement
from the definitions, but it is not a novelty claim.

\paragraph{Scope.}
Only the direct-product question in Remark 4.2 is settled here. The broader
question whether $M_d$-AP is equivalent to ordinary AP remains open; the 2025
paper restates it as the equivalence of $M_2$-AP and $M_d$-AP for all
$d\geq3$.

\begin{thebibliography}{9}
\bibitem{Vergara2023}
I. Vergara,
\emph{The $M_d$-approximation property and unitarisability},
Proc. Amer. Math. Soc. 151 (2023), 1209--1220;
\href{https://arxiv.org/abs/2202.12198}{arXiv:2202.12198}.

\bibitem{Vergara2025}
I. Vergara,
\emph{Some remarks on $M_d$-multipliers and approximation properties},
\href{https://arxiv.org/abs/2509.07861}{arXiv:2509.07861} (2025).
\end{thebibliography}

\end{document}

